% Analyse
\section{Aktuelle ITSM-Prozesse der SARIA}\label{sec:4ISTZustandsanalyse}
Im folgenden Kapitel wird der aktuelle Stand der SARIA in Bezug auf die im Zielbild genannten Prozesse untersucht. Hierbei werden konkrete Prozesse ausgewählt und diese näher beleuchtet.
Dabei wird insbesondere die deutsche IT untersucht.

\subsection{Incident \& Problem Management}\label{subsec:4incidentproblemmanagement}
In der IT-Abteilung von SARIA liegt der Schwerpunkt vornehmlich auf der Behebung von Incidents und der Abwicklung von Serviceanfragen. Grund hierfür ist, dass alle über das Ticketsystem eingehenden Anfragen, unabhängig von ihrer Natur, pauschal als Tickets klassifiziert und behandelt werden. Dieses Ticketsystem, bereitgestellt vom Schwesterunternehmen, dient als zentrale Plattform zur Erfassung und Bearbeitung der Anfragen. Die Initiation dieser Tickets erfolgt üblicherweise durch E-Mails, die von den Nutzern versendet werden.
\\
Oft wird jedoch kein Zeitfenster zur Lösung oder Bearbeitung mit den Kundinnen und Kunden vereinbart oder durch \ac{sla}s vorgegeben.
Aufgrund dieser nicht stattfindenden Kommunikation, werden die Erwartungen der Kunden an den Service der IT gesenkt, da diese davon ausgehen, dass die Bearbeitung der Tickets einige Zeit in Anspruch nehmen wird und schnelle Lösungen seltener auftreten.
\\
Die Priorisierung erfolgt oft nach dem \Gls{fifo}-Prinzip. Anderweitige Priorisierungen, wie geschäftliche Auswirkungen oder Deadlines, werden nur in Betracht gezogen, wenn diese explizit von den Kunden erwähnt werden oder Personen in höheren Positionen ausdrücklich auf eine schnelle Bearbeitung hinweisen.
Bei Tickets mit geringen Auswirkungen wird nicht explizit darauf geachtet, wenig Ressourcen zu verbrauchen. Die schnellstmögliche Lösungsfindung ist das erklärte Ziel. Sollten mehr Ressourcen benötigt werden, werden \gls{adhoc} weitere Mitarbeiter in das Thema mit einbezogen, um bei der Lösung zu unterstützen.
Sollten Tickets häufiger auftreten, wird dieses nicht ordnungsgemäß dokumentiert und es gibt keinen standardisierten Prozess, um die Lösung des Tickets schnell und effizient durchzuführen. Stattdessen wird die Erfahrung des Bearbeiters oder der Kollegen mit einbezogen, um das Incident zu lösen.
Es gibt kein zentrales Tool, wie etwa eine Wissensdatenbank oder ähnliches, in der Informationen über Tickets und deren Lösung gespeichert werden können, wie etwa \gls{known error} oder Links zu verwandten Problemlösungen oder ähnliches.
Die Möglichkeiten zur Mitteilung von Informationen, wie etwa geschäftliche Auswirkungen, betroffene User, \acs{etc}, werden nur teilweise genutzt.
Dies verschlechtert die Wahrnehmung des Serviceangebots, da so der IT-Service eher unpersönlich wirkt und ein Eindruck für die operativen Einheiten entsteht, der die IT unfähig erscheinen lässt.
\\
Eine Kategorisierung von Tickets durch den Helpdesk findet in Teilen für den Software-Bereich und die Teams statt.
\\
Dafür werden in anderen Bereichen die Tickets nur teilweise kategorisiert, welches die Bearbeitung durch die Mitarbteitenden in der IT-Abteilung erschwert.
Dies führt dazu dass Ressourcen teilweise länger als benötigt belegt werden und so die Lösung aller Incidents in angemessenem Rahmen verzögert.
\\
Ebenfalls ist eine gemeinsame und teamübergreifende Bearbeitung kaum möglich, da das Ticketsystem nur einen Bearbeiter eines Tickets vorsieht und es keine Möglichkeiten gibt, diese weiter aufzuteilen, wie etwa durch die Erstellung von Unteraufgaben.
Durch dieses Vorgehen wird die Zusammenarbeit und ebenfalls die Bearbeitung der Tickets verzögert.
\\
Mit dem Team Helpdesk \& Support gibt es zwar eine Eingangsquelle für alle Anfragen der Kundinnen und Kunden der SARIA IT, jedoch wird häufig noch das \gls{heyjo} verwendet.
Insbesondere bei Kolleginnen und Kollegen, die länger im Unternehmen arbeiten, ist der Griff zum Telefonhörer ein schneller Weg um IT-relevante Probleme zu lösen.
Dies führt jedoch dazu, dass in den normalen Ablauf der Bearbeitung von Tickets eingegriffen wird und somit weitere Verzögerungen für die ordnungsgemäß gemeldeten Tickets enstehen.

Aufgrund fehlender Unterteilung von Incidents und Service Requests und fehlenden Auswertungsmöglichkeiten im aktuellen Ticketsystem gibt es zurzeit kein Problem Management bei der SARIA IT in Deutschland.
Problems, die im Helpdesk öfter auffallen, werden teilweise angegangen, können aber nicht ordnungsgemäß nachgehalten und dokumentiert werden. Zudem fehlt für die dauerhafte Lösung oft die Zeit im Second-Level-Support, da viele Projekte stattfinden, die für die aktuellen Änderungen an der IT-Struktur und Organisation notwendig sind.

Zusammengefasst ist das Incident und Problem Management bei der deutschen IT der SARIA noch in einem frühen Stadium. Für die notwendige Professionalisierung fehlen nicht nur die Prozesse, sondern auch eine selbst anpassbare technische Lösung.

\subsection{Service Desk}
Die IT-Abteilung in Deutschland hat bereits einen Service Desk etablieren können.
Seit über 2 Jahren ist die Abteilung Helpdesk \& Support in dieser Funktion tätig.
\\
Dort werden die ankommenden Tickets erfasst und wenn nötig an die entsprechenden Teams im 2nd Level Support weitergeleitet. Auch Service Requests werden über den Service Desk beantragt und teilweise bearbeitet.
\\
Aufgrund des bisherigen Fokus auf Deutschland ist der Service Desk teilweise lokal und teilweise zentral angesiedelt. Für die Corporate Functions am Hauptstandort der SARIA steht der Service Desk auch lokal zur Verfügung mit einem \glqq{}Walk-In-Service Desk\grqq{}. Dort können Mitarbeitende von SARIA ohne Anmeldung ihre Incidents oder Service Requests melden.
\\
Weitere offizielle Eingangskanäle für den Service Desk sind das Ticketsystem und die Möglichkeit über Telefon mit dem Service Desk in Kontakt zu treten.
Inoffizielle Kanäle werden weiterhin genutzt, wie etwa WhatsApp, Anrufe und Nachrichten über Microsoft Teams. Diese unterschiedlichen Arten der Kontaktaufnahme mit dem Service Desk erschwert die Arbeit dessen, da die Informationen aus mehreren Systemen teilweise manuell in das Ticketsystem übertragen werden müssen, um eine ordnungsgemäße Dokumentation zu ermöglichen.
\subsection{Asset \& License  Management}\label{subsec:4assetlicensemanagement}
Insbesondere im Bereich des Asset \& License Management bestehen große prozessuale Lücken.
Seit ein par Jahren werden die Assets und Lizenzen in einem von einem anderen Schwesterunternehmen gehosteten System automatisiert gesammelt und vom strategischen Einkauf zur Übersicht und für die Möglichkeit eines Audits bereitgehalten.
\\
Die IT-Abteilung verwendet die gesammelten Informationen kaum, da das notwendige Know-How nicht vorhanden ist.
Zudem ist die Oberfläche des Systems unübersichtlich, da nicht nur Assets von deutschen Nutzern im System enthalten sind, sondern auch die der anderen Länder der SARIA.
Zudem wird der Support von der Seite des Schwesterunternehmens zeitnah entfallen. Dies führt zu einem weiteren Know-How Verlust, der nicht über die SARIA IT abgefangen werden kann.
\subsection{Service Catalogue \& Service Request Management}\label{subsec:4servicecatalogueandservicerequestmanagement}

Die Serviceleistungen und Anfragen an die SARIA IT sind teilweise definiert und folgen einem Standardprozess.
Kunden können Anfragen über das Ticketsystem einreichen.\footnote{siehe Kapitel \ref{subsec:4incidentproblemmanagement} auf S.\pageref{subsec:4incidentproblemmanagement}}

Ein Problem tritt auf, wenn die zuständigen Mitarbeitenden, insbesondere jene, die für die Ticketzuweisung verantwortlich sind, die Serviceanfragen nicht korrekt identifizieren können.
Dies kann zu Verzögerungen im Bearbeitungsablauf führen, falls ein Ticket irrtümlich falsch zugewiesen wird.
\\
Eine weitere Problematik besteht darin, dass den Kunden keine umfassende Übersicht über alle verfügbaren IT-Serviceanfragen zur Verfügung steht.
Dies begrenzt ihre Fähigkeit, präzise und detaillierte Anfragen an die IT der SARIA zu richten, da ihnen wesentliche Informationen fehlen.
In der Folge ist es für die bearbeitenden Mitarbeitenden oft schwierig, eine standardisierte Abwicklung der Serviceanfragen zu gewährleisten, da notwendige Daten fehlen oder der Prozess nicht einheitlich gestaltet ist und von jedem Bearbeitenden unterschiedlich gehandhabt wird.
\\
Diese Unstimmigkeiten beeinträchtigen nicht nur die Fähigkeit des Bearbeitenden, dem Kunden eine angemessene Bearbeitungszeit mitzuteilen, sondern wirken sich auch negativ auf die Wahrnehmung des IT-Services durch die Kunden aus.
Zusätzlich erschwert das Fehlen einer dokumentierten Übersicht über die Serviceanfragen und insbesondere deren Erfüllung der IT-Abteilung die Möglichkeit, einen Überblick über erbrachte Dienstleistungen und die dafür aufgewendeten Arbeitsstunden zu erhalten.
Infolgedessen ist eine Optimierung der Serviceanfragen kaum möglich, da Schwachstellen oder Prozessblockaden nicht effizient identifiziert und verbessert werden können.
Ebenso lässt sich die für die Bearbeitung einer Serviceanfrage benötigte Zeit nicht auf Grundlage früherer Anfragen bestimmen.

\subsection{Service Level Management \& Service Level Agreements}\label{subsec:4slas}
Die genannten Punkte in Kapitel \ref{subsec:4servicecatalogueandservicerequestmanagement} beeinflussen die Möglichkeit, \acs{sla} festzulegen und Kunden transparent darzustellen.
Bisherige Versuche interne Richtlinien für die Bearbeitung von Tickets festzulegen, sind aufgrund der Masse der Anfragen und der Unübersichtlichkeit des aktuellen Ticketsystems gescheitert.
\\
Mit externen Lieferanten gibt es bereits Vereinbarungen über Service Level und \ac{sla}s. Mit dem Schwesterunternehmen, das aktuell noch der Hauptlieferant ist, bestehen keine Vereinbarungen solcher Art.
Auch interne Service Level, wie etwa maximale Dauer eines Anrufs oder Tickets im First Level Support, konnten bisher nicht dauerhaft installiert werden.
Zudem findet bisher auch kaum eine Überwachung der verinbarten Service Level statt, insbesondere Berichte über die Einhaltung und Prüfungen wurden bisher kaum erstellt oder durchgeführt.
\subsection{Knowledge Management}\label{subsec:4knowledgemanagement}
Die SARIA IT in Deutschland nutzt mit der Group IT gemeinsam, seit einigen Monaten Confluence zur Erstellung von internen Dokumentationen für den IT-Betrieb und weitere Thematiken, wie Governance Richtlinien oder Anleitungen für den Helpdesk zur Vereinfachung der Bearbeitung von Tickets durch Bearbeitende.
Diese Knowledge Database ist allerdings exklusiv für die Mitarbeitenden der IT-Abteilung zugänglich.
Es gibt Ansätze für die Bereitstellung von Informationen an Kunden, wie über das Intranet und SharePoint Online.
Jedoch mangelt es an einer ganzheitlichen, kohärenten Strategie in der Implementierung und Nutzung dieser Plattformen.
Das Resultat ist eine fragmentierte Informationslandschaft, in der Benutzer sich mit unterschiedlichen Anleitungen und Informationen an verschiedenen Orten konfrontiert sehen.
\subsection{Change Management}\label{subsec:4changemanagement}

Die IT-Abteilung von SARIA verfügt aktuell über kein formal definiertes Change Management-Verfahren.
Änderungen in der IT-Infrastruktur oder bei IT-Services werden häufig \gls{adhoc}, also ohne festgelegte Prozesse oder Standards, in den Betrieb integriert.
Während in einigen Fällen eine gewisse Planung und Dokumentation von Changes erfolgt, wird dies aufgrund von Zeitdruck oft vernachlässigt.
\\
Diese Vorgehensweise birgt mehrere Risiken und Herausforderungen.
Ohne ein strukturiertes Change Management-Verfahren können unkoordinierte Änderungen zu Instabilitäten im IT-System führen, die Betriebsunterbrechungen oder Datenverlust zur Folge haben können.
Die fehlende Dokumentation erschwert die Nachverfolgung und das Verständnis für Auswirkungen dieser Änderungen, was die Fehlersuche und Problembehebung erschwert.
\\
Darüber hinaus beeinträchtigt die Ad-hoc-Implementierung von Änderungen die Fähigkeit der IT-Abteilung, die IT-Infrastruktur strategisch zu planen und zu optimieren.
Dies kann zu Ineffizienzen und zu spontanen Änderungen in den Prioritäten führen.
Zudem ist es ohne ein formalisiertes Change Management schwierig, die Einhaltung von Compliance-Vorschriften und Sicherheitsstandards sicherzustellen.

\subsection{Release Management}
Das Release Management in der SARIA wurde bisher nicht eigenständig ausgeführt, da es in der Verantwortung des Schwesterunternehmens lag.
In dieser Konstellation erhielt die SARIA IT alle relevanten Informationen zu geplanten Releases zusammengefasst und hatte nur begrenzte Möglichkeiten, eigene Planungen vorzunehmen oder Anpassungen durchzuführen.
\\
Die Zukunft sieht jedoch eine Änderung dieser Struktur vor: Die SARIA IT wird das Release Management eigenverantwortlich übernehmen müssen. 

\subsection{IT Financial Management}
Das aktuelle Financial Management der IT in der SARIA ist äußerst gering gehalten. Dies gilt insbesondere, weil die Kosten der IT am Ende des Jahres über eine Umlage mit den anderen Einheiten abgerechnet werden.
Dort wo es möglich ist werden Lizenzen direkt über die Kostenstellen der verantwortlichen Bereiche abgerechnet.
Obwohl eine jährliche Budgetplanung für IT-Ausgaben durchgeführt wird, ist die Finanzverwaltung durch zusätzliche Genehmigungsprozesse durch den Finanzvorstand im Laufe des Jahres eingeschränkt.
Diese Genehmigungsprozesse können zu Verzögerungen und möglicherweise zu Einschränkungen bei der schnellen Anpassung an sich ändernde IT-Bedarfe führen.